\documentclass[12pt,a4paper]{article}

\usepackage{fontspec}
\setmainfont{DejaVu Serif}
\setsansfont{DejaVu Sans}
\setmonofont{DejaVu Sans Mono}

\usepackage{geometry}
\geometry{left=24mm,right=24mm,top=28mm,bottom=28mm}
\usepackage{setspace}
\onehalfspacing
\usepackage{booktabs}
\usepackage{longtable}
\usepackage{array}
\usepackage{enumitem}
\usepackage{xcolor}
\usepackage{tcolorbox}
\tcbuselibrary{skins,breakable}
\usepackage{hyperref}

\definecolor{primaryblue}{RGB}{0,51,179}
\definecolor{secondaryblue}{RGB}{0,102,204}
\definecolor{successgreen}{RGB}{0,128,80}
\definecolor{warningorange}{RGB}{200,100,0}

\newcolumntype{P}[1]{>{\raggedright\arraybackslash}p{#1}}

\title{\LARGE\textbf{CodeAlchemist}\\[4pt]
\large Week 9 Codebased Addendum\\[2pt]
\normalsize Chat Bazlı Operasyon ve Hata Günlüğü}
\author{CodeAlchemist Team}
\date{Nisan 2026}

\begin{document}
\maketitle

\begin{tcolorbox}[colback=blue!5,colframe=primaryblue,title={Kapsam}]
Bu doküman, bu sohbet oturumunda Week 9 kapsamındaki token/billing geliştirmeleri
sırasında uygulanan teknik adımları, karşılaşılan hataları ve çözüm doğrulamalarını
tekil bir codebased günlüğü olarak özetler.
\end{tcolorbox}

\section{Operasyon Günlüğü}

\begin{longtable}{@{} P{0.7cm} P{4.4cm} P{4.2cm} P{4.1cm} @{} }
\toprule
\textbf{\#} & \textbf{Belirti} & \textbf{Kök Neden} & \textbf{Aksiyon ve Sonuç} \\
\midrule
\endhead
1 & Usage ledger içinde internal server error. &
\texttt{token\_package} tablosunda \texttt{description} kolonu eksikti. &
Şema migration script'i yazıldı ve çalıştırıldı;
kolon + \texttt{token\_purchase} tablosu doğrulandı. \\

2 & Token dashboard modal ekran taşması. &
Responsive kırılımlarda modal viewport'a göre sınırlanmıyordu. &
Modal konteyneri, grid kırılımı ve scroll davranışı güncellendi;
frontend build başarılı. \\

3 & Stripe configured değil uyarısı. &
Çalışan backend'in okuduğu kaydedilmiş \texttt{.env} içeriğinde
Stripe satırları eksikti. &
\texttt{server/.env} Stripe bloklarıyla güncellendi;
\texttt{/api/billing/config} yanıtında \texttt{enabled=true} doğrulandı. \\

4 & Stripe package is not installed. &
Aktif çalışma ortamında \texttt{stripe} modülü kurulu değildi. &
\texttt{.venv-win} ortamına \texttt{pip install stripe} uygulandı;
import testi geçti. \\

5 & GitHub push güvenliği sorusu. &
Secret yönetimi ve takip edilen dosyalar netleştirilmeliydi. &
\texttt{server/.env} ignore/do-not-track kontrolü yapıldı,
izlenen dosyalarda secret paterni tarandı. \\
\bottomrule
\end{longtable}

\section{Dosya Bazlı Değişiklik Özeti}

\begin{itemize}[leftmargin=*, label=\textbullet]
	\item \texttt{server/migrate\_token\_billing\_schema.py} (yeni):
	token billing şema tamamlama.
	\item \texttt{server/app.py}: env fallback ve billing config doğrulama akışı iyileştirmeleri.
	\item \texttt{client/src/components/TokenDashboardModal.jsx}:
	responsive uyum, viewport limitleri, grid kırılımları.
	\item \texttt{server/.env}: Stripe değişkenleri runtime için tamamlandı
	(git ignore kapsamında tutuldu).
\end{itemize}

\section{Doğrulama Çıktıları}

\begin{tcolorbox}[colback=green!5,colframe=successgreen,title={Başarılı Kontroller}]
\begin{itemize}[leftmargin=*, nosep]
	\item Billing config endpoint: \texttt{enabled=true}.
	\item DB şeması: \texttt{token\_package.description} mevcut.
	\item DB şeması: \texttt{token\_purchase} tablosu mevcut.
	\item \texttt{stripe} import testi başarılı.
	\item Frontend production build başarılı.
\end{itemize}
\end{tcolorbox}

\section{Risk ve Notlar}

\begin{tcolorbox}[colback=yellow!10,colframe=warningorange,title={Güvenlik Notu}]
API anahtarları sohbet içinde açık geçtiği için, repo tarafı temiz olsa bile
operasyonel güvenlik açısından key rotation gereklidir (Stripe, OpenAI,
Anthropic, Gemini, Resend).
\end{tcolorbox}

\end{document}
